\section*{ВВЕДЕНИЕ}
\addcontentsline{toc}{section}{ВВЕДЕНИЕ}

В настоящее время в виду высокой конкуренции, растущей динамики изменений предпочтений потребителей, ужесточения требований к документообороту предприятий общественного питания, каждому заведению в сфере обслуживания требуется система для эффективного управления бизнес процессами.  

Для оптимизации управления персоналом, автоматизации финансового и складского учета, налаживания логистики, повышения скорость и качества обслуживания гостей, необходимо использовать системы управления бизнесом. Они представляют собой комплексные решения, включающее в себя программное обеспечение и аппаратные средства. Программное обеспечение, способствует интеграции необходимых функций, таких как ведение финансовых отчетов, учет рабочего времени сотрудников, управление заказами, автоматизация закупок, отслеживание запасов и другие. Внедрение таких систем также позволяет минимизировать и корректировать ошибки, что помогает предприятию поддерживать конкурентоспособность.  

Раньше системы автоматизации представляли из себя узконаправленные программные пакеты, которые работали обособлено, не обмениваясь данными с другими системами. Для того чтобы система автоматизации могла удовлетворять требования каждой конкретной компании, приходилось её дорабатывать, что, в свою очередь, требовало больших финансовых и временных вложений. 

Отличие современных систем автоматизации ресторанного бизнеса заключается в том, что они объединяют все процессы в одну гибкую систему. Это открывает доступ к данным не только системе автоматизации, но и web-приложениям, приложениям для ведения складского и финансового учета. Кроме того, повысилась безопасность и конфиденциальность, а самое главное, данные системы позволяют вести бизнес процессы непрерывно и устойчиво. 

Главной задачей системы автоматизации ресторанного бизнеса является оптимизация всех бизнес-процессов в нескольких взаимосвязанных программах.   

\emph{Цель настоящей работы} – разработка программы для оптимизации ресторанного бизнеса.  Для достижения поставленных целей необходимо решить \emph{следующие задачи:}
\begin{itemize}
\item провести анализ предметной области;
\item разработать концептуальную модель приложения;
\item спроектировать приложение;
\item реализовать приложение с помощью языка программирования Python.
\end{itemize}

\emph{Структура и объем работы.} Отчет состоит из введения, X разделов основной части, заключения, списка использованных источников, X приложений. Текст выпускной квалификационной работы равен \formbytotal{lastpage}{страниц}{е}{ам}{ам}.

\emph{Во введении} сформулирована цель работы, поставлены задачи разработки, описана структура работы, приведено краткое содержание каждого из разделов.

\emph{В первом разделе} на стадии описания технической характеристики предметной области приводится сбор информации о возмржностях систем управления ресторанного бизнеса.

\emph{Во втором разделе} на стадии технического задания приводятся требования к разрабатываемому продукту.

\emph{В третьем разделе} на стадии технического проектирования представлены проектные решения для системы управления.

\emph{В четвертом разделе} приводится список классов и их методов, использованных при разработке программы, производится тестирование разработанной системы.

В заключении излагаются основные результаты работы, полученные в ходе разработки.

В приложении А представлен графический материал.
В приложении Б представлены фрагменты исходного кода. 
